\documentclass{article}
\usepackage[letterpaper, textwidth=5.5in, textheight=9.0in]{geometry}
\usepackage{setspace}
\usepackage{amsmath}

% typesetting choices
\linespread{1.08}
\frenchspacing

% math macros
\renewcommand{\vec}[1]{\mathbf{#1}}

\title{\bfseries%
The zodiacal light in {\slshape SPHEREx}:\\ A causal, geometric separation} 
\author{dwh2 }
\date{September 2025}

\begin{document}

\maketitle

\begin{abstract}\noindent
    All-sky, time-domain surveys gather intensity measurements that comprise observatory, Solar System, Milky Way, and extragalactic contributions.
    Many projects would benefit from a ``component separation'' that splits the intensity into these additive components.
    Here we perform one of these separations, under the assumption that the zodiacal light is approximately constant in a frame that rotates around the Solar System with Jupiter, and that all of the other components are approximately constant in either spacecraft coordinates or else celestial coordinates.
    Our implementation makes use of some of the ideas from the theory of invariant functions, and some ideas from causal inference.
\end{abstract}

\section{Assumptions and choices}

\begin{description}
    \item[geometry:] We assume that the zodi is approximately time invariant in a frame that rotates with Jupiter's orbit.
    In detail, our assumption is slightly weaker:
    We assume that the intensity from the zodi is an invariant scalar function of four vectors.
    The four vectors are the direction $\vec{\hat{n}}$ of the observation (or from which the photons are arriving), the barycentric displacement $\vec{r}$ of the \textsl{SPHEREx} instrument in the quasi-Newtonian Solar System frame, the barycentric displacement $\vec{r}_\text{J}$ of Jupiter in that frame, and a normal vector $\vec{\hat{z}}$ that is parallel to the total angular momentum of the zodi dust.
    This assumption is not the assumption that the zodi is rotationally invariant; it is the assumption that if you rotate all four vectors, the intensity you measure rotates in the corresponding way.
    
    In the theory of invariant functions, this further corresponds to the assumption that the zodi can be expressed as a function only of the inner products of these vectors with themselves and each other.
    Since two of the vectors are unit vectors, the only nontrivial inner products are
    \begin{align}
        x^\top = \begin{bmatrix}\vec{\hat{n}} \cdot \vec{r} &
        \vec{\hat{n}} \cdot \vec{r}_\text{J} &
        \vec{\hat{n}} \cdot \vec{\hat{z}} &
        \vec{r} \cdot \vec{r} &
        \vec{r} \cdot \vec{r}_\text{J} &
        \vec{r} \cdot \vec{\hat{z}} &
        \vec{r}_\text{J} \cdot \vec{r}_\text{J} &
        \vec{r}_\text{J} \cdot \vec{\hat{z}} \end{bmatrix} ~,
    \end{align}
    where $x$ is a column vector of explicitly scalar features, which encode everything possible about a time-invariant zodi that depends only on Jupiter, the Sun, and the \textsl{SPHEREx} point of view.

    There will be two other components of the intensity field.
    One of these is the spacecraft component.
    This will be a scalar function of the direction vector $\vec{\hat{n}}$, the barycentric displacement $\vec{r}$ of the spacecraft, a spacecraft pole vector HOGG and the displacement HOGG WHAT.

    The other component is the combined astronomical component, which absorbs all the stars, interstellar medium, galaxies, quasars, and everything else outside the Solar System.
    

    HOGG WHAT ARE WE GIVING UP WITH THESE ASSUMPTIONS?

    \item[non-negativity:] foo
\end{description}

\end{document}
